%%=============================================================================
%% Samenvatting
%%=============================================================================

% TODO: De "abstract" of samenvatting is een kernachtige (~ 1 blz. voor een
% thesis) synthese van het document.
%
% Een goede abstract biedt een kernachtig antwoord op volgende vragen:
%
% 1. Waarover gaat de bachelorproef?
% 2. Waarom heb je er over geschreven?
% 3. Hoe heb je het onderzoek uitgevoerd?
% 4. Wat waren de resultaten? Wat blijkt uit je onderzoek?
% 5. Wat betekenen je resultaten? Wat is de relevantie voor het werkveld?
%
% Daarom bestaat een abstract uit volgende componenten:
%
% - inleiding + kaderen thema
% - probleemstelling
% - (centrale) onderzoeksvraag
% - onderzoeksdoelstelling
% - methodologie
% - resultaten (beperk tot de belangrijkste, relevant voor de onderzoeksvraag)
% - conclusies, aanbevelingen, beperkingen
%
% LET OP! Een samenvatting is GEEN voorwoord!

%%---------- Nederlandse samenvatting -----------------------------------------
%
% TODO: Als je je bachelorproef in het Engels schrijft, moet je eerst een
% Nederlandse samenvatting invoegen. Haal daarvoor onderstaande code uit
% commentaar.
% Wie zijn bachelorproef in het Nederlands schrijft, kan dit negeren, de inhoud
% wordt niet in het document ingevoegd.

\IfLanguageName{english}{%
    \selectlanguage{dutch}
    \chapter*{Samenvatting}
    \lipsum[1-4]
    \selectlanguage{english}
}{}

%%---------- Samenvatting -----------------------------------------------------
% De samenvatting in de hoofdtaal van het document

\chapter*{\IfLanguageName{dutch}{Samenvatting}{Abstract}}

% \begin{document}

% \begin{abstract}

Deze bachelorproef onderzoekt hoe een AI-gedreven agent kan worden ingezet om dependency-updates uit bestaande tools voor automatisch dependencybeheer te analyseren en te beheren, zodat het kernteam dat verantwoordelijk is voor platform- en dependencybeheer deze updates efficiënter en met minder handmatig werk kan verwerken. Het achterliggende probleem is dat bestaande dependencybots updates wel detecteren, maar de inhoudelijke interpretatie van changelogs en release notes grotendeels aan ontwikkelaars overlaten.

Om deze onderzoeksvraag te beantwoorden, werd een vergelijkende studie uitgevoerd op basis van een proof-of-concept genaamd DepGuard AI. Binnen dit prototype werden drie AI-varianten uitgewerkt en onder dezelfde omstandigheden geëvalueerd: een Mastra gebaseerde agent, een n8n workflow en een eigen OpenAI agent met toolloop. De evaluatie gebeurde aan de hand van vijf representatieve testcases en een LLM-as-a-Judge benchmark, waarbij de output werd beoordeeld op specificiteit, volledigheid, actiegerichtheid en verwerkingstijd.

De resultaten tonen aan dat alle drie de varianten bruikbare samenvattingen kunnen genereren, maar dat er duidelijke verschillen bestaan in kwaliteit en snelheid. Mastra behaalde de hoogste inhoudelijke scores met een overallscore van 7,81, tegenover 6,53 voor OpenAI en 6,04 voor n8n. OpenAI was de snelste variant, terwijl n8n het laagst scoorde op specificiteit en volledigheid. De statistische analyse bevestigt dat de verschillen tussen de agents op meerdere scorecategorieën betekenisvol zijn.

Op basis van de implementatie en evaluatie kan worden besloten dat een AI-gestuurde benadering een relevante meerwaarde biedt voor dependencybeheer, omdat ze de interpretatielast van changelogs en release notes verlaagt. Tegelijk blijkt dat de kwaliteit van de uitkomst sterk afhangt van de gekozen orchestratie en van de beschikbaarheid van actuele broninformatie. Voor de context van Springbok vormt Mastra daardoor de meest geschikte basis voor verdere uitwerking van het dependency-dashboard.

% \end{abstract}

% \end{document}

% \lipsum[1-4]
